%%%% ABSTRACT
%%
%% Versão do resumo para idioma de divulgação internacional.

\begin{abstractutfpr}%% Ambiente abstractutfpr
This study investigates the effects of fine-tuning on language models tasked with navigating ASCII mazes through the use of the LLM-MRI library. By analyzing the activation graphs generated from the models' embeddings, we assess structural changes using various graph metrics with distribution inference via statistical sampling. Our findings indicate subtle, yet not significant, alterations in the graph structures after fine-tuning, suggesting that, although fine-tuning influences information processing, these changes are not drastically reflected in the analyzed metrics. The study highlights the importance of considering different levels of abstraction within the models to fully understand the impact of fine-tuning. Recommended future work includes exploring diverse prompts, using embbeddings that capture task resolution, and examining the effects of various fine-tuning techniques and model architectures on activation graph structures.
\end{abstractutfpr}
